\scanpage{128}
\uline{Définition} Soit $K/Q$ une extension \deleted{\unsure{algébrique}} \inserted{algébrique}. On dit que $K$ est tot. réel si pour tout $Q$-hom $\sigma:K\to C$, on a $\sigma(K)\subset\unsure{\mathbb R|R}$. On dit qu'un élément $\unsure{x|a}$ de $K$ est tot. réel \inserted{(resp. tot. $>0$, tot. $\ge0$)} si pour les $\sigma$ comme dessus, on a $\sigma(x)\in\unsure{\mathbb R|R}$ (resp. $\sigma(x)>0$, resp. $\sigma(x)\ge0$). Tot. (ou purement) imaginaire: \ldots.

\uline{Sorite} Pour que $x$ soit tot. réel, il \unsure{f. et s.} que $Q(x)\subset K$ le soit. Pour qu'il soit tot. $\ge0$, il faut et il suffit qu'il soit nul ou tot $>0$. Les éléments tot $\ge0$ \deleted{\unsure{forment}} \inserted{sont} les éléments $\ge0$ pour une relation d'ordre sur $K$. [On parlera de $x$ tot. $\ge y$ ou $x$ tot. $\le y$]. Si $K$ est engendré par des ss-corps $K_\alpha$, alors $K$ est tot. réel ssi les $K_\alpha$ le sont. Si \inserted{$K$ est engendré par des $x_\alpha$,} $K$ tot. réel ssi \deleted{\unsure{les élts}} les $x_\alpha$ le sont.

Il existe, pour $K$ donné, un plus grand ss-corps totalement réel de $K$, appelé la \uline{clôture tot. réelle} de $Q$ dans $K$. Soit $\deleted{K=}\overline Q\subset C$, et $\pi=\operatorname{Gal}(\unsure{\overline Q}/Q)$, enfin $\tau\in\pi$ la conjugaison complexe
\[
\tau\lambda=\overline\lambda\qquad\text{pour }\lambda\in\overline Q
\]
de sorte qu'on a
\[
\tau^2=1.
\]
Alors le ss-groupe $\pi'$ de $\pi$ qui \deleted{\unsure{invariante}} \inserted{correspond} à la clôture tot. réelle $\deleted{\unsure{de}}\mathbb K\deleted{\unsure{Q}}$ de $Q$ dans $\overline Q$ est \deleted{\unsure{le}} le ss-groupe \uline{invariant} \uline{fermé} engendré par $\tau$, i.e. le ss groupe invariant engendré par les ${}^{\sigma}\tau=\sigma\tau\sigma^{-1}$ $(\sigma\in\pi)$.
\margin{\diagram{\begin{tikzcd}[row sep=small]
A \arrow[d,dash] \\
\mathbb L \arrow[d,dash] \\
\mathbb K \arrow[d,dash] \\
Q
\end{tikzcd}}}
Considérons la plus grande extension abélienne $A$ de $\mathbb K$, dont le groupe de Galois est donc $H=\pi'/(\pi',\pi')$. \deleted{\unsure{En}} Sur
\mspage{127}
